% 3.2.1 跨场景部署图；示例具体设备来自H02 R1/H的存档，未绘运行时间。
\begin{figure}[htbp]
 \centering
 \resizebox{0.98\textwidth}{!}{%
 \begin{tikzpicture}[>=Latex,font=\small,
 stageA/.style={draw=blue!65!black,fill=blue!8,rounded corners=2pt,align=center,minimum width=1.65cm,minimum height=.74cm},
 stageB/.style={draw=orange!70!black,fill=orange!9,rounded corners=2pt,align=center,minimum width=1.65cm,minimum height=.74cm},
 stageC/.style={draw=green!50!black,fill=green!9,rounded corners=2pt,align=center,minimum width=1.65cm,minimum height=.74cm}]
   \draw[rounded corners=5pt,draw=gray!60] (-.10,-2.25) rectangle (6.75,2.27);
   \draw[rounded corners=5pt,draw=gray!60] (8.05,-2.25) rectangle (14.90,2.27);
   \node[font=\bfseries] at (3.33,1.95) {场景1：应用A与B};
   \node[font=\bfseries] at (11.45,1.95) {场景2：应用A与C};
   \node[font=\scriptsize,anchor=east] at (.45,.78) {A};
   \node[font=\scriptsize,anchor=east] at (.45,-1.13) {B};
   \node[font=\scriptsize,anchor=east] at (8.59,.78) {A};
   \node[font=\scriptsize,anchor=east] at (8.59,-1.13) {C};
   \node[stageA] (al1) at (1.5,.78) {A.pre\\CPU2};
   \node[stageA] (al2) at (3.42,.78) {A.main\\NPU0};
   \node[stageA] (al3) at (5.34,.78) {A.post\\CPU1};
   \node[stageA] (ar1) at (9.68,.78) {A.pre\\CPU2};
   \node[stageA] (ar2) at (11.60,.78) {A.main\\NPU0};
   \node[stageA] (ar3) at (13.52,.78) {A.post\\CPU1};
   \node[stageB] (bl1) at (1.5,-1.13) {B.pre\\CPU1};
   \node[stageB] (bl2) at (3.42,-1.13) {B.main\\NPU0};
   \node[stageB] (bl3) at (5.34,-1.13) {B.post\\CPU1};
   \node[stageC] (cr1) at (9.68,-1.13) {C.pre\\CPU1};
   \node[stageC] (cr2) at (11.60,-1.13) {C.main\\GPU0};
   \node[stageC] (cr3) at (13.52,-1.13) {C.post\\CPU1};
   \foreach \a/\b in {al1/al2,al2/al3,ar1/ar2,ar2/ar3,bl1/bl2,bl2/bl3,cr1/cr2,cr2/cr3}
     {\draw[-Latex,thick] (\a)--(\b);}
   \draw[dashed,blue!75!black,thick,rounded corners=3pt] (.51,.23) rectangle (6.32,1.35);
   \draw[dashed,blue!75!black,thick,rounded corners=3pt] (8.68,.23) rectangle (14.50,1.35);
   \draw[<->,blue!75!black,thick,dashed] (6.75,.78) -- (8.05,.78);
   \node[align=center,font=\scriptsize,blue!70!black] at (7.40,1.47) {共享\\具体设备};
   \node[font=\scriptsize] at (3.42,-1.84) {各场景独立选择OPP与启动时刻};
   \node[font=\scriptsize] at (11.60,-1.84) {各场景独立选择OPP与启动时刻};
 \end{tikzpicture}}
 \caption{跨场景共享Stage执行单元的示意（设备映射取自H02的R1/H方案，省略具体OPP和时间）}
 \label{fig:ch3-shared-scene}
\end{figure}
